Рассмотрим язык
\[
\textsf{SPECIAL-HAMPATH} = \left\{(G, s, t, K) \colon
\begin{aligned}
    &G = (V, E) \text{ -- ориентированный граф}, K \subset E,\\
    &\forall X \subset K\text{ существует ориентированный}\\
    &\text{гамильтонов путь из $s$ в $t$, проходящий через все}\\
    &\text{ребра из $X$ и не проходящий через все ребра из $K \setminus X$}
\end{aligned}
   \right\}.
\]

Во-первых, нужно доказать принадлежность к $\Pi_2^p$, а во-вторых -- её $\Pi_2^p$-трудность. 

1. $\textsf{SPECIAL-HAMPATH} \in \Pi_2^p$.

Достаточно расшифровать описание языка и перевести его на язык кванторов:
\begin{multline*}
    (G, s, t, K) \in \textsf{SPECIAL-HAMPATH} \Longleftrightarrow\\\Longleftrightarrow \forall \,X\;\exists \,\text{путь} \,v_1, \ldots, v_n\colon 
\left\{
\begin{aligned}
    &X \subset K, \; v_1, \ldots, v_n \text{ -- гамильтонов},\;v_1=s,\\
    &v_n=t,\;\text{всякое ребро из $X$ лежит на пути,}\\
    &\text{всякое ребро из $K \setminus X$ не лежит на пути.}
\end{aligned}
\right\}
\end{multline*}

Все условия в фигурных скобках проверяются за полиномиальное время, а переменные, по которым берётся квантор, имеют полиномиальный размер относительно размера входа. Значит, если в качестве верификатора $V(x, y, z)$ взять функцию, которая по входу $x = (G, s, t, K)$, подмножеству вершин графа $y = X$ и пути в графе $z = (v_1, \ldots, v_n)$ проверяет корректность данных и условия в фигурных скобках, то для него выполнено определение принадлежности к $\Pi_2^p$.

2. $\textsf{SPECIAL-HAMPATH}$ -- $\Pi_2^p$-трудный.

Как известно, язык
$
\Pi_2^p\textsf{-3SAT} = \{\phi\text{ в форме 3-КНФ}\mid \forall\,\mathbf{y}_1\;\exists\, \mathbf{y}_2\colon \phi(\mathbf{y}_1, \mathbf{y}_2)=1\}
$ является $\Pi_2^p$-трудным, поэтому достаточно свести его к нашему языку.

Сделаем это аналогично доказательству сводимости $\textsf{3SAT} \le_p \textsf{HAMPATH}$, а именно построим по формуле $\phi$, зависящей от двух наборов переменных $\mathbf{y}_1$ и $\mathbf{y}_2$ ориентированный граф с <<гаджетами>> для каждой переменной, как это было на лекции. В качестве $K$ можно взять верхние левые рёбра <<гаджетов>>, соответствующих переменным из $\mathbf{y}_1$. Как мы знаем, проход по такому ребру (то есть <<зигзагом>>) означает истинность переменной, а противоположная ситуация, когда по такому ребру не совершается обход (то есть <<загзигом>>), означает ложность переменной.

Таким образом, выбор значений для набора $\mathbf{y}_1$ равносилен выбору некоторых рёбер $X \subset K$ (выбрали ребро --- соответствующая переменная истинна, иначе --- ложна), а выбор значений для $\mathbf{y}_2$ равносилен нахождению гамильтонового пути в построенном графе, который каким-то образом (зигзагом или загзигом) пройдёт через гаджеты для $\mathbf{y}_2$ и пройдёт через гаджеты для $\mathbf{y}_1$ в соответствии с выбором множества $X$.

В качестве $\Sigma_3^p$-полного языка возьмём следующий искусственный язык, являющийся компиляцией задачи о рюкзаке и трёхшаговой игры:
\[
\textsf{3BACKPACK} = \left\{
\left(
\begin{aligned}
&n, k, m, N,\\
&A = \{a_i\colon i = 1, \ldots, n\},\\
&BC = \{(b_i, c_i)\colon i = 1, \ldots, k\},\\
&D = \{d_i\colon i = 1, \ldots, m\}
\end{aligned}
\right)
\colon
\begin{aligned}
    &\text{можно взять несколько $a_i$ так, чтобы}\\
    &\text{при любом выборе по одному элементу из}\\
    &\text{каждой пары $(b_j, c_j)$ полученный набор}\\
    &\text{можно дополнить некоторыми $d_l$ до суммы $N$}
\end{aligned}
\right\}
\]

Сначала осознаем определение сего языка: первый игрок выбирает подмножество из $A$, второй игрок из каждой пары множества $BC$ выбирает по одному элементу, а затем первый игрок выбирает ещё несколько элементов, но уже из $D$. Первый игрок выигрывает, если сумма всех чисел равна $N$. Язык состоит из тех ситуаций, когда у первого игрока есть выигрышная стратегия.

Принадлежность к $\Sigma_3^p$ почти очевидна: 
\[
    (n, k, m, N, A, BC, D)\in\textsf{3BACKPACK} \Longleftrightarrow  \exists {X \subset A}\;\forall {Y \subset BC}\;\exists {Z \subset D}\colon\;\; \ldots\ldots\ldots,
\]
где условие на месте многоточия проверяется за полином: корректность выбора чисел и проверка на равенство их суммы $N$. 

Сведём $\Sigma_3^p$-полный язык $\Sigma_3\textsf{3SAT} = \{\phi\text{ в форме 3-КНФ}\mid \exists\,\mathbf{y}_1\;\forall\, \mathbf{y}_2\;\exists\,\mathbf{y}_3\colon \phi(\mathbf{y}_1, \mathbf{y}_2, \mathbf{y}_3)=1\}$ к нашему. Решение почти дословно повторяет решение задачи 12 с \href{https://youtu.be/dN753JpeSDs?t=3645}{консультации Ильи Степанова}. Пусть в формуле $\phi$ всего $T$ дизъюнктов и $S$ переменных. Для каждого литерала (то есть для $y$ и $\neg y$) заведём своё битовое число $\mathcal{N}_y$ и $\mathcal{N}_{\neg y}$, состоящее из $3T+2S$ битов. Первые $3T$ будут для дизъюнктов -- по 3 бита на каждый блок. Запишем в $i$-ый блок 001, если соответствующий числу литерал используется в $i$-ом дизъюнкте, и 000 иначе. Оставшиеся $2S$ битов поделим на $S$ блоков по два бита. Для литералов $y$ и $\neg y$, где $y$ -- $j$-ая переменная, запишем в $j$-ый блок числа, соответствующего этому литералу, 01, а во все остальные -- 00. Также для каждого дизъюнкта с номером $i$ возьмём по два числа со всеми нулями, за исключением блока, отвечающего $i$-ому дизъюнкту -- там напишем 001. Число $N$ положим равным (опять же, в двоичной записи)
\[
N = \underbrace{011\,011\,\ldots\,011}_{\text{$T$ блоков}}\;\underbrace{01\,01\,\ldots\,01}_{\text{$S$ блоков}}.
\]

В качестве $A$ возьмём числа  $\{\mathcal{N}_{y}, \mathcal{N}_{\neg y}\colon y \in \mathbf{y}_1\}$, в качестве $BC$ -- множество пар чисел $\{(\mathcal{N}_{y}, \mathcal{N}_{\neg y})\colon y \in \mathbf{y}_2\}$, а в качестве $D$ --- все остальные числа (то есть множество $\{\mathcal{N}_{y}, \mathcal{N}_{\neg y}\colon y \in \mathbf{y}_3\}$ и специальные числа для дизъюнктов).

Выбор чисел соответствует тому, чему мы полагаем равным переменные в формуле -- истине или лжи. В начале игры если первый игрок полагает $y = 1$, то он из $A$ берёт число $\mathcal{N}_y$, иначе -- $\mathcal{N}_{\neg y}$, и так по всем переменным из $\mathbf{y}_1$. Второй игрок обязан выбрать ровно одно число из пары, что соответствует его выбору значений переменных из $\mathbf{y}_2$. Первому игроку остаётся выбрать значения из $\mathbf{y}_3$ и взять несколько вспомогательных чисел из $D$ -- надо это для того, чтобы в каждом блоке из первой половины чисел в сумме было 011: если формула выполнима, то в сумме там уже хотя бы 001, так что остаётся лишь дополнить это число до требуемого с помощью вспомогательных чисел.

%\begin{note}
%Определим языки $\Sigma_k 3SAT$ и $\Sigma_k 3DNFSAT$ аналогично $\Sigma_k SAT$ для формул в 3-КНФ и 3-ДНФ соответственно. Аналогично определим $\Pi_k 3SAT$ и $\Pi_k 3DNFSAT$. Тогда $\Sigma_k 3SAT$ полон в $\Sigma^p_k$ при нечётных k. $\Pi_k 3SAT$ полон в $\Pi^p_k$ при чётных k. Это достигается обычным сведением формулы к виду 3-КНФ за счёт того, что последний квантор — квантор существования, который и навешивается на новые переменные.
%\end{note}
% зачем это добавлено?

